\section{Simulation environment}
\label{sec:environment}

The environment is the rig's digital twin with a contact-manipulable
payload added. This section covers where the twin comes from, how the
payload and its contacts are designed, why the timing is structured
around an exact substep count, and what the policy senses.

\paragraph{Twin provenance.}
One parametric CAD description is the single source of truth for the
physical rig: the printed parts are exported from it, and the simulation
model --- meshes, joint frames, arm-base placements, camera poses --- is
regenerated from the same description. Simulated geometry therefore
cannot drift from the physical build, and any rig change propagates to
the twin by regeneration rather than by hand-editing. The three mounted
cameras (one scene view, one wrist view per arm) sit at the poses the
printed mounts dictate.

\paragraph{Payload.}
The payload is a cube of \SI{2.2}{\centi\metre} edge and
\SI{40}{\gram} mass resting on the table. A cube suits both tasks with
one object: its symmetric cross-section reads the same to either
gripper, and a parallel-jaw grasp of a cube tolerates yaw error that a
long prism does not (Section~\ref{sec:difficulties} records how the
earlier prism payload motivated this). \unjustified{} the mass and the
exact edge length were chosen by convention (a comfortable fraction of
the jaw opening) and have not been validated against a physical object.

\paragraph{Contact design.}
Simulated parallel-jaw grasping fails in characteristic ways: objects
slip through soft contacts, or explode out of stiff ones. We use the
standard anti-slip recipe: an elliptic friction cone with a high
friction-to-normal impedance ratio, extra no-slip solver iterations, a
raised contact priority on small high-friction fingertip pads fitted to
both jaws, and a physics timestep of $1/600$\,\si{\second}. The gripper
actuators are force-bounded so that a closed-jaw command against a
blocked jaw saturates instead of building unbounded penetration force.
\unjustified{} the friction coefficients, impedance ratio and pad
dimensions were tuned by iterating on rendered episodes until grasps
held through transport; they are stable in practice but not derived
from measured material properties.

\paragraph{Timing.}
Control, dataset frames and evaluation all run at \SI{30}{\hertz}. With
the $1/600$\,\si{\second} physics step, one control tick is exactly twenty
substeps, so a simulated episode advances in ticks that are bit-identical
in duration to training frames. The finer step also stabilises the
pinch contacts. Evaluation rollouts therefore observe the world on
exactly the clock the policy was trained on.

\paragraph{Sensing and goal specification.}
The policy observes the three camera views and a twelve-dimensional
proprioceptive state: five joint angles plus a gripper opening fraction
per arm, in the same units, ordering and semantics as the real rig's
recorder, so a policy trained in simulation reads real observations
without translation. The task goal is communicated visually: a
translucent disc on the table marks the target zone in the scene
camera's view. No goal coordinates are ever given to the policy.

\paragraph{Gripper opening cap.}
The real rig's trigger maps to a deliberately capped jaw opening
(roughly a third of the mechanical range), a garment-handling choice
described in the companion paper~\citep{so101teleop}. The simulated
gripper honours the same cap on both the command side and the recorded
open fraction, so simulated and real gripper channels share semantics;
the cube still passes the capped opening with clearance.
